%==============================================================================
%  CAPÍTULO XV
%==============================================================================
\chapter{Manual de tramitación: reorganización judicial ordinaria}
\label{cap:tram-reorg-ordinaria}

\fichaProcedimiento
  {Empresa Deudora mediana o grande (sobre \uf{25.000} de ventas anuales).}
  {Juzgado civil del domicilio del deudor.}
  {De seis a doce meses hasta el acuerdo firme, según la complejidad del pasivo.}
  {Alto: honorarios del veedor, asesoría legal y financiera, y costos de auditoría.}
  {Acuerdo de reorganización vinculante para todos los acreedores valistas, o apertura
   de liquidación si se rechaza.}

El desarrollo dogmático de este procedimiento se encuentra en el capítulo
\ref{cap:reorganizacion-ordinaria}. Este capítulo aborda exclusivamente su ejecución
forense.

%==============================================================================
\bloqueEncargo
%==============================================================================

La reorganización ordinaria es el procedimiento más exigente del sistema y el que con
mayor frecuencia se acepta sin condiciones de éxito. Antes de asumir el encargo deben
concurrir tres presupuestos materiales, ninguno de los cuales es jurídico:

\begin{enumerate}
  \item \textbf{Un negocio viable.} La reorganización reestructura pasivos; no genera
  ingresos. Si el flujo operacional proyectado no permite servir el pasivo
  reestructurado, el acuerdo se incumplirá y el procedimiento habrá consumido recursos
  para llegar igualmente a la liquidación.

  \item \textbf{Capacidad de producir información financiera.} El procedimiento exige
  estados financieros, proyecciones y respuesta rápida a los requerimientos del veedor.
  Una empresa sin contabilidad al día no puede sostenerlo.

  \item \textbf{Un mapa de mayorías alcanzable.} Si los dos o tres acreedores principales
  concentran el pasivo y son hostiles, el resultado está determinado antes de presentar.
\end{enumerate}

\begin{foco}[La conversación previa con los acreedores mayoritarios]
La reorganización se gana fuera del tribunal. Antes de presentar debe sondearse la
disposición de los acreedores que concentran el pasivo con derecho a voto. Una propuesta
presentada sin ese sondeo es una apuesta, no una estrategia.
\end{foco}

%==============================================================================
\bloqueEntrevista
%==============================================================================

Sobre el cuestionario general del capítulo \ref{cap:entrevista}, este procedimiento
requiere indagar además lo siguiente.

\begin{cuestionario}[Viabilidad del negocio]
  \pregunta{¿Cuál es el flujo operacional mensual de los últimos doce meses, y cuál el
  proyectado para los próximos veinticuatro?}
  {Es el insumo central de la propuesta. Sin él no puede diseñarse un plan de
  amortización realista.}

  \pregunta{¿La crisis proviene de una caída de ingresos, de un aumento de costos, de un
  cliente que dejó de pagar o de una decisión de financiamiento?}
  {Determina si la causa es reversible. Una crisis por pérdida del cliente principal no
  se resuelve reestructurando plazos.}

  \pregunta{¿Qué contratos son esenciales para la continuidad del giro?}
  {La resolución de reorganización mantiene vigentes los contratos de suministro. Hay que
  identificarlos antes para invocar el efecto.}

  \pregunta{¿Hay bienes que puedan enajenarse sin afectar la operación?}
  {Permite diseñar una propuesta con componente de realización de activos no esenciales,
  que mejora sustancialmente su atractivo para los acreedores.}
\end{cuestionario}

\begin{cuestionario}[Mapa de acreedores]
  \pregunta{¿Qué proporción del pasivo concentran los tres mayores acreedores?}
  {Define si el acuerdo depende de negociación bilateral o de convencer a una masa
  dispersa.}

  \pregunta{¿Alguno de los acreedores es persona relacionada con el deudor o sus socios?}
  {Su exclusión de determinados cómputos altera el quórum efectivo, y su participación
  puede fundar impugnaciones por abuso del derecho.}

  \pregunta{¿Cuál ha sido la actitud de cada acreedor principal hasta ahora: negociación,
  cobranza, ejecución?}
  {Anticipa quién apoyará el acuerdo y quién intentará forzar la liquidación.}

  \pregunta{¿Existen acreedores con garantía real sobre bienes esenciales?}
  {Su posición es privilegiada y su oposición puede hacer inviable la continuidad
  operativa aun con acuerdo aprobado.}
\end{cuestionario}

%==============================================================================
\bloqueDocumental
%==============================================================================

\begin{checklist}[Antecedentes específicos de la solicitud]
  \item Declaración jurada de activos y pasivos conforme al \art{56}.
  \item \textbf{Informe de auditoría externa} (art. 56): obligatorio para la gran empresa
  deudora, emitido por un auditor registrado ante la \cmf{} que certifique la exactitud del estado
  de situación financiera. Sólo la micro y pequeña empresa que opta por la reorganización
  simplificada de la \Lreforma{} queda exceptuada.
  \item Nómina de acreedores con individualización, monto, naturaleza del crédito,
  preferencias y garantías.
  \item Nómina de juicios pendientes con individualización de tribunal y rol.
  \item Certificado de deudas previsionales y laborales.
  \item Nómina de bienes con indicación de gravámenes y de su carácter esencial o no
  esencial para el giro.
  \item Proyección de flujos que sustente la propuesta de acuerdo.
  \item Propuesta de terna de veedores.
\end{checklist}

%==============================================================================
\bloqueFocos
%==============================================================================

\subsection{La declaración jurada del artículo 56}

Es el documento de mayor riesgo del procedimiento. Su inexactitud no genera sólo una
objeción procesal: alimenta impugnaciones del acuerdo y, si el procedimiento deriva en
liquidación, puede constituir el antecedente de la causal del \art{169 A} \Nro 1.
Debe construirse con los antecedentes verificados del capítulo \ref{cap:entrevista}, no
con la contabilidad interna sin contraste.

\subsection{El aprovechamiento del período de protección}

La \pfc{} es un recurso escaso y no renovable indefinidamente. El error de gestión más
frecuente consiste en usar las primeras semanas para ordenar información que debió estar
lista antes de presentar, y llegar a la junta con una propuesta improvisada.

\begin{foco}[Planificación regresiva del período de protección]
La propuesta de acuerdo debe estar sustancialmente redactada \emph{antes} de presentar la
solicitud. El período de protección se usa para negociar y ajustar, no para redactar.
\end{foco}

\subsection{El informe del veedor}

El veedor no es un adversario, pero tampoco un aliado: su función es informar a los
acreedores sobre la verosimilitud de la propuesta. La calidad de la relación de trabajo
con él ---entrega oportuna de información, respuestas completas, ausencia de sorpresas---
influye directamente en el tono del informe, y el informe influye en el voto.

\subsection{El financiamiento durante la protección}

El \art{74} otorga preferencia a los créditos contratados durante el período de
protección. Es una herramienta poco utilizada y decisiva: permite obtener capital de
trabajo cuando el crédito ordinario está cerrado. Debe explorarse desde el primer día.

%==============================================================================
\bloqueCronograma
%==============================================================================

\begin{checklist}[Secuencia de trabajo]
  \item \textbf{Semanas $-8$ a $-4$.} Diagnóstico, verificación documental y construcción
  del mapa de acreedores.
  \item \textbf{Semanas $-4$ a $-1$.} Sondeo de los acreedores mayoritarios. Redacción de
  la propuesta y de la proyección de flujos. Preparación de la terna de veedores.
  \item \textbf{Día 0.} Presentación de la solicitud.
  \item \textbf{Resolución de reorganización.} Comienza el período de protección
  financiera. Publicación en el \boletin{}.
  \item \textbf{Período de verificación.} Control diario de las verificaciones que
  ingresan; objeción oportuna de las improcedentes.
  \item \textbf{Antes de la junta.} Informe del veedor. Ajuste final de la propuesta y
  cierre de la negociación con los acreedores decisivos.
  \item \textbf{Junta de acreedores.} Deliberación y votación.
  \item \textbf{Post-junta.} Resolución que tiene por aprobado el acuerdo; eventual
  incidente de impugnación; publicación y entrada en vigor.
\end{checklist}

\begin{alerta}[El control diario del Boletín Concursal]
Desde la resolución de reorganización, la notificación se practica por el \boletin{}.
Los plazos corren aunque nadie avise. La revisión diaria de la causa no es una buena
práctica: es el estándar mínimo de diligencia.
\end{alerta}

%==============================================================================
\bloqueErrores
%==============================================================================

\errorfrecuente{Presentar sin haber sondeado a los acreedores mayoritarios}
{La propuesta se rechaza en junta y el procedimiento deriva en liquidación, habiendo
consumido meses y honorarios.}
{Sondear informalmente antes de presentar y ajustar la propuesta a lo que sea votable.}

\errorfrecuente{Usar el período de protección para preparar la propuesta}
{Se llega a la junta con un plan sin sustento financiero, que el informe del veedor
expondrá.}
{Llegar al día cero con la propuesta y la proyección de flujos ya elaboradas.}

\errorfrecuente{Entregar información al veedor de forma reactiva y parcial}
{El informe recoge las reservas del veedor y contamina el voto de los acreedores.}
{Establecer un canal ordenado desde la nominación, con entregas completas y anticipadas.}

\errorfrecuente{Ignorar el financiamiento preferente del artículo 74}
{La empresa se queda sin capital de trabajo durante la protección y llega deteriorada a
la junta.}
{Explorar el crédito preferente desde la primera semana del período.}

\errorfrecuente{Diseñar la propuesta pensando sólo en el deudor}
{Los acreedores no tienen incentivo para aprobar un plan que les ofrece menos de lo que
obtendrían en liquidación.}
{Construir la propuesta contra un escenario de liquidación comparado, y mostrarlo.}
